\section{HIỆN TRẠNG ỨNG DỤNG CNTT TẠI ĐƠN VỊ}
\subsection{ Hiện trạng ứng dụng CNTT trong quản lý tài sản đô thị}
Hiện nay, công tác quản lý tài sản đô thị tại TP. Đà Nẵng còn phân tán, sử dụng nhiều công cụ khác nhau mà chưa được tích hợp trên một nền tảng thống nhất:
\begin{table}[htpb]
    \centering
    \caption{Khảo sát hiện trạng các hệ thống quản lý hạ tầng hiện có}
    \label{tab:hien_trang_he_thong}
    \renewcommand{\arraystretch}{1.5} % Tăng khoảng cách các hàng để bảng thoáng và dễ đọc hơn
    \begin{tabular}{|c|p{3.5cm}|c|c|p{2.5cm}|p{2.4cm}|}
        \hline
        \textbf{STT} & 
        \centering\textbf{Tên hệ thống/ứng dụng} & 
        \textbf{Môi trường} & 
        \textbf{Tích hợp \gls{gis}} & 
        \centering\textbf{Ngôn ngữ lập trình và hệ Quản trị \gls{csdl}} & 
        \centering\arraybackslash\textbf{Đơn vị phát triển} \\ \hline
        
        1 & Phần mềm quản lý cây xanh đô thị & Desktop & Không & Bảo mật & Sở Xây dựng/2018 \\ \hline
        2 & Hệ thống quản lý đèn chiếu sáng & Desktop & Không & Bảo mật & Công ty TNHH/2019 \\ \hline
        3 & Phần mềm quản lý đường bộ & Desktop & Không & Bảo mật & Sở GTVT/2016 \\ \hline
        4 & Hệ thống quản lý thoát nước & Desktop & Có (cơ bản) & Bảo mật & Trung tâm Thoát nước/2020 \\ \hline
        5 & Phần mềm quản lý biển báo GT & Desktop & Không & Bảo mật & Sở GTVT/2017 \\ \hline
    \end{tabular}
\end{table}
\subsection{Hiện trạng hạ tầng CNTT}\
Hạ tầng \gls{cntt} hiện tại của nhóm bao gồm:
\begin{table}[htpb]
    \centering
    \caption{Hạ tầng CNTT hiện tại của nhóm}
    \label{tab:hien_trang_he_thong}
    \renewcommand{\arraystretch}{1.5} % Tăng khoảng cách các hàng để bảng thoáng và dễ đọc hơn
    \begin{tabular}{|c|c|c|}
        \hline
        \textbf{STT} & 
        \centering\textbf{Thiết bị/Hệ thống} & 
        \textbf{Số lượng}  
        
             \\ \hline
        
        1 & Laptop & 3  \\ \hline
        2 & Thiết bị phát sóng wifi  & 3  \\ \hline
        
    \end{tabular}
\end{table}
\subsection{Hiện trạng về nhân lực \gls{cntt}}
Nhân lực \gls{cntt} hiện có của nhóm:
\begin{itemize}
    \item Cán bộ chuyên trách \gls{cntt}: 15 người
    \item Cán bộ quản lý dự án \gls{cntt}: 3 người
    \item Chuyên viên bảo mật thông tin: 3 người
\end{itemize}

Về tổ chức bộ máy đơn vị có các nhân viên phụ trách \gls{cntt} như sau: Công nghệ thông tin gồm 18 nhân viên (3 PM, 15 \gls{ai}); trong đó, có 05 nhân viên chuyên trách làm công tác cơ yếu, 05 nhân viên chuyên trách làm công tác \gls{cntt}. Đến nay, còn 05 nhân viên chuyên trách làm công tác \gls{cntt}. Trình độ chuyên môn \gls{cntt}, Cơ yếu, chuyên ngành khác, trình độ đại học. 

\subsection{Sự cần thiết phải đầu tư}
Hiện trạng quản lý tài sản đô thị tại TP. Đà Nẵng đang đối mặt với nhiều thách thức nghiêm trọng:
\begin{itemize}
    \item Dữ liệu tài sản đô thị phân tán, không có \gls{csdl} tập trung và thống nhất, gây khó khăn trong việc tra cứu, tổng hợp báo cáo.
    \item Thiếu công cụ trực quan hóa trên bản đồ, dẫn đến việc lập kế hoạch bảo trì hạ tầng kém hiệu quả.
    \item Không có cơ chế cảnh báo sớm về hư hỏng, xuống cấp của tài sản, dẫn đến phản ứng chậm và tốn kém chi phí sửa chữa.
    \item Quy trình tiếp nhận và xử lý phản ánh của công dân về sự cố hạ tầng còn thủ công, thiếu minh bạch.
    \item Không có công cụ phân tích dự báo nhu cầu bảo trì, thay thế tài sản dựa trên dữ liệu lịch sử.
\end{itemize}

\subsection{Mối liên hệ của dự án với hệ thống ứng dụng CNTT khác}
Dự án \gls{geoai} được xây dựng với định hướng mở, đảm bảo khả năng liên thông và tích hợp chặt chẽ với các hệ thống hiện hữu của thành phố:
\begin{itemize}
    \item \textbf{Cổng Dữ liệu mở Thành phố / Overture Maps:} Kết nối \gls{api} để tự động đồng bộ hóa và làm giàu dữ liệu tài sản, địa điểm tiện ích (POI), giao thông từ cơ sở dữ liệu mở toàn cầu Overture Maps và OpenStreetMap.
    \item \textbf{Hệ thống Phản ánh kiến nghị công dân (Tổng đài 1022):} Chia sẻ dữ liệu báo cáo sự cố hạ tầng thông qua \gls{api}, tự động phân loại mức độ rủi ro và xác định tọa độ GPS của sự cố trên bản đồ \gls{gis} tập trung.
    \item \textbf{Cơ sở dữ liệu Không gian dùng chung (PostGIS):} Hệ thống chia sẻ trực tiếp các lớp bản đồ không gian (Spatial Layers) cho các sở ngành khác (Sở Xây dựng, Sở GTVT) thông qua kiến trúc dữ liệu tập trung, giúp xóa bỏ tình trạng phân mảnh dữ liệu.
\end{itemize}

\cleardoublepage